\section{The actual semilocal metric and its angle cutoffs}

The source for the semilocal Hilbert space in this section is
C99, VII.(5), its Fourier unitarity is VII, Lemma 1(b), and
its product-norm multiplication cutoff is VII.(12)--(13).
The real prolate replacement is VIII.(31)--(33).
The standard finite character vector is VIII.(14), whose
finite Fourier transform in the rational setting was computed
in AC-T11.  The distinction between those source definitions
and the new fixed-angle comparison below is explicit.
Throughout, the raw norm is the positive-half-line norm
\(\|\psi\|^2=\int_0^\infty|\psi(u)|^2\,du\), with the
specified parity extension to \(\mathbb R\).
This fixes the factor \(1/2\) in the tensor representative;
the unrescaled full-line norm would have an additional
factor two.
Here \(\chi\) labels the Dirichlet \(L\)-factor and
the finite vector \(\varphi_\chi\) of AC-T11.
C99's character of the compact-unit action is its
inverse in these coordinates; consequently the
source's \(\overline{\chi_v}\) in VIII.(14) is
precisely the ramified unit vector used here.

\begin{theorem}[AC-T47: the exact rational semilocal operator dictionary]
\label{AC-T47}
Let \(\chi\) be a primitive Dirichlet character of conductor
\(q\), including the principal character \(q=1\), and let
\(S\) contain infinity and every prime dividing \(q\).
Write \(\mathcal P=\{p\in S:p\nmid q,\ p<\infty\}\) and
\(M_S=\prod_{p\in\mathcal P}p\).
On the raw parity space put
\[
 D_r\psi(u)=r^{-1/2}\psi(u/r),\qquad
 Z_{\chi,S}
    =\prod_{p\in\mathcal P}
        (1-\chi(p)p^{-1/2}D_{1/p})^{-1}.
 \tag{AD.1}
\]
Empty products are the identity.
The compact-character sector of the actual semilocal
Hilbert space is isometrically represented by the raw
space with inner product
\[
           (\psi,\phi)_{\chi,S}
                        =\langle Z_{\chi,S}\psi,
                                   Z_{\chi,S}\phi\rangle_{L^2(du)}.
 \tag{AD.2}
\]
The normalization uses one-half the even/odd tensor
representative, so that the two rational signs are
counted once.
In these coordinates its orthogonal product-norm
time cutoff and its Fourier transformation are
\[
 A_{\chi,S,\Lambda}=Z_{\chi,S}^{-1}P_\Lambda Z_{\chi,S},
 \qquad
 {\cal F}_\chi=\omega_\chi D_q\mathcal F_\epsilon,
 \qquad |\omega_\chi|=1,
 \tag{AD.3}
\]
where \(P_\Lambda=1_{|u|\leq\Lambda}\) and
\(\chi(-1)=(-1)^\epsilon\).
The second map is a unitary map from the metric
\((\cdot,\cdot)_{\chi,S}\) to
\((\cdot,\cdot)_{\bar\chi,S}\).
The orthogonal frequency cutoff is therefore
\[
 B_{\chi,S,\Lambda}
       ={\cal F}_\chi^{-1}A_{\bar\chi,S,\Lambda}{\cal F}_\chi .
 \tag{AD.4}
\]
Its range, as an actual set of raw functions, is
the ordinary Fourier band \(|\xi|\leq\Lambda/q\).
The range of \(A_{\chi,S,\Lambda}\) is the ordinary
raw time interval \(|u|\leq\Lambda\).

For the minimal set \(S=\{\infty\}\cup\{p:p\mid q\}\),
\(Z_{\chi,S}=1\).  Thus the actual semilocal pair is
exactly the ordinary time-band pair with radii
\((\Lambda,\Lambda/q)\), in the ordinary raw \(L^2\)
metric.  Its full-index prolate family is the family
of AC-T41, or AC-T42--AC-T45 when \(q=1\).
For additional unramified places, (AD.3) is
orthogonal in (AD.2), not in the unweighted raw metric.
\end{theorem}
\begin{proof}
Let \(\varphi_{\chi,S}\) be the finite standard vector:
on a ramified prime it is the unit-supported character
vector of AC-T11, and on an unramified prime in \(S\)
it is \(1_{\mathbb Z_p}\).
Use the representative
\(\frac12\psi\otimes\varphi_{\chi,S}\).
A fundamental domain for the \(S\)-unit action has
positive real coordinate and unit finite coordinates.
The compact unit character has modulus one there.
In the orbit sum of VII.(5), a ramified coordinate
forces the valuation of the rational \(S\)-unit
to be zero.  At an unramified coordinate its
valuation must be nonnegative.
After the two signs have been combined using the
specified parity, the normalized orbit-sum function is
\[
 \sqrt{x}\sum_{\substack{n\geq1\\
                   \text{all prime factors of }n\text{ in }\mathcal P}}
                     \chi(n)\psi(nx)
                       =(UZ_{\chi,S}\psi)(\log x),
 \quad (U\psi)(t)=e^{t/2}\psi(e^t).
 \tag{AD.5}
\]
The definition of the quotient norm in VII.(5)
therefore gives (AD.2), with compact-unit Haar
measure one.
Both the series and its inverse are bounded:
\[
 \begin{split}
 \|Z_{\chi,S}\|&\leq
                  C_S^+:=\prod_{p\in\mathcal P}(1-p^{-1/2})^{-1},\\
 \|Z_{\chi,S}^{-1}\|&\leq
                  C_S^-:=\prod_{p\in\mathcal P}(1+p^{-1/2}).
 \end{split}
 \tag{AD.6}
\]
Indeed the dilation operators are unitary and commute,
the direct series has absolutely summable coefficients,
and the inverse is the finite divisor polynomial.
Consequently the standard tensor images are dense
in the entire compact-character sector, not merely
in a proper cyclic subspace.  This also proves the
claimed isometric coordinate identification.

On that fundamental domain, product norm is exactly
the positive real coordinate.  Source VII.(12)
is multiplication by \(1_{x\leq\Lambda}\) on (AD.5).
Transport by the just proved unitary identification
gives the first formula in (AD.3).
All dilations in \(Z\) and \(Z^{-1}\) move raw
support to the left, so each maps the time-supported
subspace into itself.  Since they are inverse,
the range of \(A\) is exactly that subspace.
Thus larger individual real coordinates on some
finite-place fibers do not give a larger time
range after passing through the quotient map.

For the Fourier formula, the finite Fourier
calculation (T.22) is
\(\widehat\varphi_\chi(y)
 =(\tau(\chi)/q)\varphi_{\bar\chi}(qy)\),
with the evident principal convention \(q=1\).
The primes dividing \(q\) belong to \(S\), so
multiplication by the rational number \(q\)
is an \(S\)-unit relation in the quotient.
Moving this factor from the finite coordinates
to the real coordinate changes the real vector
to \((\tau(\chi)/\sqrt q)D_q\mathcal F_\epsilon\psi\).
This proves (AD.3), with its harmless unit phase.
It is also possible to verify unitarity directly:
the Mellin multiplier of \(Z_{\chi,S}\) is
\[
 L_{\mathcal P}(s,\chi)
       =\prod_{p\in\mathcal P}(1-\chi(p)p^{-s})^{-1}.
 \tag{AD.7}
\]
On \(\operatorname{Re}s=1/2\), its modulus for
\(\bar\chi\) at \(\tau\) equals its modulus for
\(\chi\) at \(-\tau\).
Raw Fourier transform is Mellin reflection
times the unimodular factor (TU.13), and
\(D_q\) contributes a unit-modulus Mellin factor.
This proves the metric assertion without
replacing an equivalent norm by an isometric one.
Formula (AD.4) follows from Fourier unitarity.
Since \(D_q\) multiplies time support by \(q\),
its raw band range is exactly
\(|\xi|\leq\Lambda/q\).

When \(\mathcal P\) is empty, the orbit sum (AD.5)
has only \(n=1\); this proves the minimal-set
assertion.  For example, with one extra unramified
prime and \(a=\chi(p)p^{-1/2}\), the transported
time projection in logarithmic raw coordinates
is exactly
\[
 Z^{-1}P_LZ
   =P_L+
       1_{(L-\log p,L]}\sum_{k\geq1}a^kT_{k\log p},
 \qquad (T_bf)(t)=f(t+b).
 \tag{AD.8}
\]
This follows from \(T_bP_L=P_{L-b}T_b\).
Taking a nonzero input supported just above \(L\)
shows that (AD.8) is generally different from
the unweighted raw cutoff.
Its orthogonality is exactly the one asserted
in (AD.2), so this example does not contradict
the source's orthogonal multiplication cutoff.
\end{proof}

\begin{theorem}[AC-T48: angle cutoffs for a fixed semilocal set have the same tested global trace]
\label{AC-T48}
Keep a fixed \(S,\chi\) as in AC-T47 and choose
a fixed concentration threshold \(0<\eta<1\).
In the metric (AD.2), let
\[
 {\cal R}_{\chi,S,\Lambda}^{\eta}
  =\operatorname{ran}
       1_{[\eta,1]}
        \bigl(A_{\chi,S,\Lambda}B_{\chi,S,\Lambda}
                         A_{\chi,S,\Lambda}
               \bigm|_{\operatorname{ran}A_{\chi,S,\Lambda}}\bigr).
 \tag{AD.9}
\]
This is a finite-dimensional space defined by the
actual semilocal pair, with its actual orthogonal
metric.  Map it to the global sector by the standard
finite tensor map: its global theta image is
\(E_\chi{\cal R}_{\chi,S,\Lambda}^{\eta}\) for
\(\chi\ne1\); for \(\chi=1\) use the full \(J\)
completion of AC-T45, and use its two-constraint
subspace when taking the positive defect.

Assume \(\mathsf A\).  The orthogonal projections
onto these global finite images have, for every
original compact norm test, the same trace
asymptotics as AC-T41 and AC-T44--AC-T45.
In particular their tested trace difference from
the original full-index prolate family tends to zero.
The assertion holds for every fixed \(S,\eta,\chi\);
it does not assert estimates uniform when \(S\)
grows with \(\Lambda\).
The angle threshold (AD.9) is explicitly a second
quantization convention.  It is not asserted to
have the same finite rank or the same finite range
as the source's real full-index cutoff.
\end{theorem}
\begin{proof}
We verify the precise analytic properties of
(AD.9), and then apply the already proved finite-image
trace argument with those verified properties.
Write \(R=\Lambda/q\), \(c=\Lambda R\), and
\(\kappa_S=C_S^+C_S^-\).

\emph{1. Compactness and a controlled band extension.}
Transport the semilocal Fourier map by the two
isometries \(Z_\chi,Z_{\bar\chi}\).
It becomes
\[
       \widetilde{\cal F}_\chi
             =Z_{\bar\chi}{\cal F}_\chi Z_\chi^{-1}
 \tag{AD.10}
\]
on ordinary raw \(L^2\).
Expand the first factor by its absolutely convergent
series and the last factor by its finite divisor
polynomial.  Using
\(\mathcal F_\epsilon D_{1/d}=D_d\mathcal F_\epsilon\),
(AD.10) is an operator-norm convergent sum of
unit phases times
\[
        \frac{\mu(d)\chi(d)\bar\chi(n)}{\sqrt{dn}}\,
                     D_{qd/n}\mathcal F_\epsilon ,
 \quad d\mid M_S,\quad n\text{ supported on }\mathcal P.
 \tag{AD.11}
\]
Every \(P_\Lambda D_r\mathcal F_\epsilon P_\Lambda\)
is Hilbert--Schmidt, by its bounded sine or cosine
kernel on a compact square.
The series of their operator norms is summable.
Thus \(P_\Lambda\widetilde{\cal F}_\chi P_\Lambda\)
is compact.  Its squared singular-value operator
is the isometric transport of the concentration
operator in (AD.9).
Consequently (AD.9) is finite dimensional.

For \(\psi\) in this range, invert the concentration
operator on its selected range to obtain \(v\).
Let \(w=B_{\chi,S,\Lambda}v\).
Then
\[
             \psi=A_{\chi,S,\Lambda}w,\qquad
             \operatorname{supp}\widehat w
                        \subset[-R,R],\qquad
             \|w\|_2\leq\kappa_S\eta^{-1}\|\psi\|_2.
 \tag{AD.12}
\]
All inverses and projections used for the first
two assertions are in (AD.2); (AD.6) gives the
last raw norm bound.
Expanding the finite inverse in \(A=Z^{-1}PZ\)
also shows
\[
             (Aw)(u)=w(u)\quad
                         \text{for }|u|<\Lambda/M_S.
 \tag{AD.13}
\]
Indeed on that interval every shifted cutoff
\(P_{\Lambda/d}\), \(d\mid M_S\), equals one,
and the finite inverse cancels \(Z\).
In particular these selected vectors are smooth
near raw zero, and their value there agrees
with the band extension's value.

\emph{2. Fractional domain bounds and actual approximation.}
For every \(0<s<1/2\),
\[
       \|\psi\|_{H_D^s}
               \leq C_{S,s,\eta,\chi}(1+c)^s\|\psi\|_2,
                  \qquad \psi\in{\cal R}_{\chi,S,\Lambda}^{\eta}.
 \tag{AD.14}
\]
To prove it, expand \(A w\) using (AD.1):
\[
 Aw=\sum_{d\mid M_S}\sum_n
       \frac{\mu(d)\chi(dn)}{\sqrt{dn}}\,
                  P_{\Lambda/d}D_{1/(dn)}w .
 \tag{AD.15}
\]
For each term, the raw Fourier bandwidth is
\(dnR\), its time radius is \(\Lambda/d\), and
their product is \(nc\).
The sharp-endpoint estimate (TB.33), applied to
its logarithmic representative, bounds its
\(H_D^s\) norm by
\(C_s(1+nc)^s\|w\|_2\).
Here the interior log derivative is at most
\((1/2+2\pi nc)\|w\|_2\), by the band derivative
bound; this proves the needed hypothesis of
(TB.33) for each term.
The series converges in \(H_D^s\), because
\[
       \sum_n n^{s-1/2}
            =\prod_{p\in\mathcal P}(1-p^{s-1/2})^{-1}<\infty .
 \tag{AD.16}
\]
Together with the finite divisor sum and (AD.12),
this proves (AD.14).
No ordinary \(H^1\) estimate is being asserted
across the additional internal cutoff boundaries.

Let \(\Pi^\eta\) denote the orthogonal projection
onto (AD.9) in (AD.2).
For every raw \(f\),
\[
 \|(1-\Pi^\eta)f\|_2
   \leq C_{S,\eta}\bigl(
           \|(1-P_\Lambda)f\|_2+
           \|(1-Q_R)f\|_2\bigr),
 \quad Q_R=\mathcal F_\epsilon^{-1}
                         1_{[-R,R]}\mathcal F_\epsilon .
 \tag{AD.17}
\]
For a vector in \(\operatorname{ran}A\), the
spectral theorem gives
\[
 \|(1-\Pi^\eta)f\|_{\chi,S}^2
       \leq(1-\eta)^{-1}\|(1-B)f\|_{\chi,S}^2 .
 \]
For a general \(f\), first separate its orthogonal
\(A\)-component.  Distance to \(\operatorname{ran}A\)
in the weighted norm is at most \(C_S^+\) times
its ordinary distance to the same raw time range.
Distance to \(\operatorname{ran}B\) has the
identical bound for its raw band range.
The triangle inequality and (AD.6) give (AD.17).
This uses the actual orthogonal projections in
the transported metric, not equality of the
weighted and unweighted projections.

For the smooth bulk \(P^p_\Lambda Q_R^pP^p_\Lambda T\),
its raw time tail is zero and its raw Fourier
tail has every polynomial rate by (TU.3).
For a smooth endpoint profile the same assertion
uses its corresponding Schwartz tail.
Combining (AD.17), (AD.14) with \(s=3/8\), and
the bound \(E_\chi:H_D^{5/16}\to L^2\), gives
exactly the graph interpolation used in AC-T38
and AC-S29.  In particular an \(X_m\) endpoint
column has graph error
\[
          C_{S,\eta,\chi,m}\,
                        c^{-m/6+5/16}\|\Psi\|_{2,X_m}.
 \tag{AD.18}
\]
Hilbert--Schmidt versions follow by summing the
raw squared errors and the coarse Sobolev bounds,
and then applying Holder to those sums.
Thus the actual selected family has all the
bulk and endpoint approximation properties
used in AC-T39 and AC-T41.
In particular (AD.14), with any
\(5/16<s<1/2\), places its finite-dimensional
range in the domain of the global closed
theta operator.  The standard tensor map
on these possibly nonsmooth vectors means
that graph closure of its Schwartz-core map;
no bounded map from the entire semilocal
completion into the global completion is
being assumed.

In the principal sector it also has the exact
two-constraint approximation property of AC-T42.
For mean errors use its raw time support and
Cauchy--Schwarz.  For value errors use (AD.13):
on \((-\Lambda/M_S,\Lambda/M_S)\) the selected
vector has band derivative bound
\(2\pi R\kappa_S\eta^{-1}\|\psi\|_2\).
The local inequality (TP.8) therefore gives
the same normalized \(c^{-m/2}\) value error,
with constants depending on \(S,\eta\).
Project the two profiles (TP.6) by \(\Pi^\eta\);
their normalized constraint matrix tends to
the identity.  The exact correction (TP.9)
then applies.  This proves the required principal
moment conditions, rather than assuming they
are preserved by \(\Pi^\eta\).

\emph{3. The global arithmetic Gram remains subpower.}
Equip the selected raw range with an ordinary
raw orthonormal basis, and form its global
theta Gram matrix \(G\), using the two-constraint
range in the principal case.
Under \(\mathsf A\), for every \(\varepsilon>0\),
\[
                       \|G^{-1}\|\leq C_{S,\eta,\chi,\varepsilon}
                                                    c^\varepsilon .
 \tag{AD.19}
\]
Here are the changes to the proof of AC-T28,
or step 1 of AC-T44.
By (AD.12)--(AD.13) and the band \(L^\infty\)
bound, a sufficiently small
\(a=C_{S,\eta}/R>0\) satisfies
\[
       \int_0^a|\psi(u)|^2du\leq\tfrac12\|\psi\|_2^2
 \tag{AD.20}
\]
for all large \(\Lambda\).
The interval lies within \((0,\Lambda/M_S)\);
the constant is chosen using
\(\|w\|_2\leq\kappa_S\eta^{-1}\|\psi\|_2\).
Thus the weighted raw norm in (TB.37)
is at least \(C_{S,\eta,\delta}c^{-\delta}\|\psi\|_2\).
Multiplication by
\(\min((u/\Lambda)^\delta,1)\) is uniformly
bounded on \(H_D^\alpha\), \(0<\alpha<1\):
in log coordinates it is bounded and Lipschitz,
and its fractional difference quotient gives
this bound directly.  Hence (AD.14) with
\(\alpha=3/8\) supplies the upper Sobolev
estimate (TB.38), in place of an ordinary
global derivative estimate.
The zero-free reciprocal bound (TB.32), or
(TP.22) in the principal case, is unchanged.
Interpolation in (TB.39)--(TB.40), with
arbitrarily small \(\delta\) and reciprocal
exponent, proves (AD.19).
The Mellin identity on the shifted line is
legitimate by graph approximation on the
time-supported core.  Its theta image has
upper support \(\Lambda\), so multiplication
by \((x/\Lambda)^\delta\) is bounded there.
In the principal sector preserve \(I=0\)
while making that core approximation.

\emph{4. The actual global theta tails are small.}
Let \(T_g\) be the raw lower-theta-tail operator
from AC-T29, or (TP.23) in the principal case.
For either character type, its outside-time
kernel proves the slightly more general estimate
\[
       \|T_g1_{u>b}\|_2^2\leq C_{\chi,g}/(\Lambda b).
 \tag{AD.21}
\]
This is exactly the integral in (TB.47) or
(TP.26), with lower raw limit \(b\).
Use the exact identity
\[
                  1-A=Z^{-1}(1-P_\Lambda)Z .
 \tag{AD.22}
\]
For a term \(d^{-1/2}D_{1/d}\) of \(Z^{-1}\),
the shifted tail starts at \(b=\Lambda/d\).
Equation (AD.21) bounds its Hilbert--Schmidt
norm by \(C_{\chi,g}\sqrt d/\Lambda\);
the coefficient cancels \(\sqrt d\).
There are only finitely many divisors of \(M_S\).
Multiplying on the right by \(Z\) gives
\[
                         \|T_g(1-A)\|_2
                                      \leq C_{S,\chi,g}\Lambda^{-1}.
 \tag{AD.23}
\]
For \(\psi=Aw\) in (AD.12), the band theta
image vanishes below \(\Lambda^{-1}\).
In the principal case this assertion requires
\(w(0)=0\), and follows exactly from
\(\psi(0)=0\) and (AD.13).
Thus \(T_g\psi=-T_g(1-A)w\).
Summing over the raw orthonormal selected
basis, using the bounded extension map in
(AD.12), and then (AD.19), proves the
normalized tail estimate
\[
 \|V(g)(1-S_\Lambda)E G^{-1/2}\|_2^2
                   \leq C_{S,\eta,\chi,g,\varepsilon}
                                  \Lambda^{-2+2\varepsilon}.
 \tag{AD.24}
\]
The upper tail is zero by raw time support,
with mean zero in the principal case.
This proves the exact tested comparison with
the positive window contraction by the
operator calculation AC-T30.

The zero-tail row kernels (TB.57)--(TB.59),
or (TP.30)--(TP.31), satisfy the same outside
estimate (AD.21).  They annihilate the band
extension, with the same principal value-zero
qualification.  Hence (AD.22)--(AD.24) apply
to every fixed zero jet as well.
There is no need to assert bounded variation
of a selected vector across all its internal
boundaries.  The identity equating its finite
window zero-jet integral to the negative
regularized tail functional is an identity
of bounded raw functionals on the time-supported
space: the window has only finitely many
sampling terms, and the tail kernels above
are bounded in raw \(L^2\).
It holds on smooth compact raw cores by the
proved Mellin identities, and extends by
raw norm density; preserve the continuous
mean condition in the principal case.
Thus the full-multiplicity lower bound
AC-T31 or (TP.32) applies to this family.

\emph{5. The upper bound, frequency tails, and comparison.}
The local finite-pole inverse \(T_L\), its exact
zero-jet cancellation, its two Fourier boundary
descriptions, and the source-independent identities
(TU.17) and (TP.35) use ordinary smooth raw
cutoffs.  None of these constructions is changed.
Their approximation by the selected finite
family is now supplied by (AD.18), with the
principal two-constraint correction proved
in step 2.
The half-line image approximations AC-T41-halfline
and AC-T43 are also unchanged: they are statements
about global \(E_\chi\) and actual compact raw
profiles.  Equations (AD.18) and (AD.19)
therefore give every finite-image step of
AC-T39 and AC-T44's local upper argument.
This proves compact-frequency equality for
the actual image projection.

For completeness the quantitative part is
preserved as well.  At a unit frequency window
of height \(T\), AC-T40 gives polynomial-height
boundary norms and a bulk norm at most
\((2+|T|)^C(1+L)^{C\log(2+|T|)+C}\).
The additional constants in (AD.18)--(AD.24)
depend on the fixed \(S,\eta,\chi\), and do
not depend on \(T,L\).
The half-line approximation still has cost
\(\exp(C B^{102})\), while (AD.18) with
\(m=12\) still has exponent \(-27/16\).
Thus a polynomial threshold
\(L\geq C_{S,\eta,\chi}(2+|T|)^D\)
absorbs all boundary and bulk errors, exactly
as in (TT.27)--(TT.30) or (TP.42)--(TP.44).
Below that threshold use the elementary
positive density bound \(L/\pi\).
The uniform polynomial unit-interval bound
and Schwartz tightness follow.
Together with (AD.24), this proves the
same original-test trace formula.

In the principal case the fixed-profile
approximation proved in step 2 is unconditional.
It gives strong convergence of the constrained
image projections by the cyclic-density proof
of AC-T45.  Its unconstrained version gives
convergence of both fixed profile moments,
using the local derivative estimate following
(AD.18).  The rank-two graph and trace-norm
argument (PC.9)--(PC.11) consequently recovers
exactly the same two polar characters.
This proves the full principal assertion.
Finally the original full-index family has
already been proved to have these same trace
limits.  Subtracting the two proved formulas
establishes the asserted tested-trace comparison.
This subtraction compares two explicitly
defined quantizations; it does not alter either
threshold or infer equality of their ranges.
\end{proof}

\begin{remark}[What is now identified, and the order of limits]
For the rational minimal semilocal set, the actual
product-norm pair and metric are exactly the real
pair, with the conductor and parity specified in
AC-T47.  Thus the full-index tensor family used
above is the direct rational extension of C99's
VIII.(31)--(33), with a now proved all-character
analytic trace under \(\mathsf A\).
For each additional fixed finite set of unramified
places, AC-T48 proves equality of the tested
global trace for its explicitly defined genuine
angle cutoff.
The constants (AD.6) and (AD.16) show why no
uniform growing-\(S\) limit has been asserted.
Taking closed global time-image ranges before
quantization is yet another operation: the
half-line density statements can erase the
finite spectral defect under that order of
completion.  It is not the finite tensor
construction of VIII.(33).
\end{remark}

\paragraph{Geometric scope.}
These are operator identities in the semilocal quotient
Hilbert space specified by C99, VII.(5), followed by its
global tensor map VIII.(33).
Their new content is the exact transported metric and
the convergence of actual finite projection traces.
They do not construct a scheme, a finite-etale cover,
or an \(\ell\)-adic realization, and no SGA theorem is
an input to their proofs.
The two principal lines retain the characters recorded
in C99, III.(31)--(34), as checked in AC-T45.
Any interpretation by an additional Grothendieck
cohomological comparison remains the distinct
dictionary question already described after AC-T11.
